{\Large Звук}
{
\definecolor{HumanAudioColor}{rgb}{0.2,0.2,0.6}
%
\begin{itemize}

\item Хоча звук, океанські хвилі та серцебиття не є електромагнітними, вони включені в цю діаграму як частотний орієнтир. Інші властивості електромагнітних хвиль відрізняються від звукових хвиль.

\item Звукові хвилі спричинені коливальним стисненням молекул. Звук не може поширюватися у вакуумі, наприклад у відкритому космосі.

\item Швидкість звуку в повітрі на рівні моря становить 1240 км/год (770 миль/год), а у воді значно більша --- 5328 км/год (3311 миль/год).

% From https://en.wikipedia.org/wiki/Hearing_range
\item Людина може чути звук лише в діапазоні від $\approx$20Hz до $\approx$20kHz. Деякі кажани можуть чути звук до 200kHz.

\item Інфразвук (нижче 20Hz) може сприйматися внутрішніми органами та дотиком.
%
%from http://www.worldtravelers.org/motionsickness.htm
Частоти в діапазоні 0,2Hz часто є причиною морської хвороби.


\item 88 фортепіанних клавіш рівномірно темперованого строю точно розташовані на частотній діаграмі. Музична нота ля (A4 або A440 при 440Hz) зображена на діаграмі як \pscirclebox[fillstyle=solid,fillcolor=HumanAudioColor]{\textcolor{white}{A4}}.

\item Протягом віків люди прагнули розділити безперервний звуковий частотний спектр на окремі музичні ноти, що мають гармонійні співвідношення, крім октави. Мікротональні музиканти вивчають різні гами. За останніми підрахунками, існує 4700 різних музичних гам.

% \textcolor{gray}{\hrule}
% \vspace{.1in}

%\item This image depicts air being compressed as sound waves in a tube from a speaker and then travelling through the tube towards the ear.\\
%%picture of sound
%\vspace{1.75in}
%\rput[t]{0}(1,-.5){\input{pictures/soundwave.tex}}
\end{itemize}
% 	\vspace{0.1in}
% 	%
% 	\begin{tabular}{rl}
% 	%\begin{minipage}[b]{2.45in}\rput(.65,1){\input{pictures/soundwave.tex}}
% 	\begin{minipage}[b]{2.45in}\rput(1.3,.55){\includegraphics[width=2.2in]{pictures/soundbmp.eps}}
% 	\end{minipage}&
% 	\hspace{0.0in}
% 	\begin{minipage}[b]{1.2in}\textcolor{white}{%
% 	This image depicts air being compressed as sound waves in a tube from a speaker and then travelling through the tube towards the ear.\\	}
% 	\end{minipage}
% 	\end{tabular}
% 	\vspace{.2in}
}

